\section{Robustness and an exploratory horizon extension}
\subsection{Overlapping events, news and influential observations}

Table~\ref{tab:robustness} shows how the later addition means respond to separate exclusions. Removing overlapping MSCI changes reduces the 2017--2026 mean from 2.15\% to 1.46\%, with an interval of [0.50\%, 2.47\%]. The latest-period mean falls from 1.50\% to 0.73\%, with an interval including zero. Overlap events have unusually large average returns: 11.25\% among the 12 later-decade overlaps and 14.16\% among the five latest-block overlaps. This makes the overlap consequential for the magnitude of the OMXSB association.

\begin{table}[!htbp]
\centering
\caption{Sensitivity of addition returns}\label{tab:robustness}
\small
\setlength{\tabcolsep}{5pt}
\begin{tabular}{llrrl}
\toprule
Period & Specification & $N$ & CAR & 95\% interval \\
\midrule
2017--26 & Baseline & 169 & 2.15 & [1.01, 3.39] \\
2017--26 & Without MSCI overlap & 157 & 1.46 & [0.50, 2.47] \\
2017--26 & Without screened company news & 163 & 1.91 & [1.01, 2.82] \\
2017--26 & Winsorized 1st and 99th percentiles & 169 & 1.92 & [0.91, 3.05] \\
2017--26 & Without Sivers & 168 & 1.71 & [0.57, 2.96] \\
2022--26 & Baseline & 87 & 1.50 & [-0.10, 3.22] \\
2022--26 & Without MSCI overlap & 82 & 0.73 & [-0.77, 2.32] \\
2022--26 & Without screened company news & 84 & 1.30 & [0.12, 2.43] \\
2022--26 & Winsorized 1st and 99th percentiles & 87 & 1.04 & [-0.11, 2.25] \\
2022--26 & Without Sivers & 86 & 0.62 & [-0.85, 2.23] \\
\bottomrule
\end{tabular}
\begin{minipage}{\linewidth}\footnotesize\vspace{5pt}
Event-weighted CAR $[AD-1,ED+1]$ in percent, with pointwise restricted wild cluster bootstrap intervals by review. Each exclusion is applied separately. The company-news screen examines the 24 most extreme inherited events, rather than the full event universe, and is outcome-informed. MSCI exclusion changes the target population. Winsorization changes the outcome. These are sensitivity descriptions, not independent replications or specification selection rules. Sources: effect, msci\_overlap, news\_screen and extreme\_values aggregate CSV files.
\end{minipage}
\end{table}


The exclusion changes the target population. It does not identify the incremental contribution of each index when both change. Overlapping events can differ in size, visibility, investor demand and information. The smaller non-overlap estimate is thus evidence of sensitivity to a recognizable subgroup, not a causal decomposition into OMXSB and MSCI effects. A clean attribution would require a design that distinguishes the assignment and timing of the two index changes.

The targeted news exclusion leaves a later-decade mean of 1.91\% and a latest-block mean of 1.30\%. These values should not be labeled news-free returns. Only the 24 most extreme inherited events were reviewed under the screen, so unexamined ordinary observations may also contain firm news. Conditioning the audit on extreme outcomes can be useful for detecting obvious problems, but it is unsuitable as a comprehensive rule for declaring all remaining returns attributable to membership.

Removing Sivers reduces the later-decade mean to 1.71\%, while the latest-block mean falls to 0.62\%. The former remains positive with a pointwise interval excluding zero; the latter does not. Winsorization gives the same qualitative distinction between a more stable pooled later result and a sensitive final-period estimate. Because exclusions and winsorization answer different questions, their agreement cannot be counted as independent evidence. The full-sample result remains the baseline, and the range of estimates describes its dependence on data handling.

These checks imply two useful constraints on the argument. First, the positive later-decade addition pattern is not entirely generated by one extreme stock or by the identified overlapping index events. Second, the latest five-year average is not sufficiently stable to support a confident claim about the current size of a membership premium. Neither statement licenses an unconditional claim about the population of all omitted historical events.

\subsection{Dependence, weighting and repeated securities}

The later addition mean of 2.15\% becomes 2.50\% when reviews receive equal rather than event-count weights. Removing one review at a time produces event-weighted means between 1.87\% and 2.47\%. The analytical CR1 standard error is 0.55 percentage points, the review jackknife standard error is also approximately 0.55, and the review-by-security two-way estimate is 0.47. In this sample, these alternatives do not reveal a hidden collapse of precision behind the later pooled mean. They still rely on assumptions about dependence and should not be interpreted as proofs that the effective independent sample size is known.

For the latest addition block, the event and equal-review means are both approximately 1.50\%. The leave-one-review-out range is 1.02\% to 2.01\%, and analytical standard errors range from about 0.71 to 0.77 across the reported estimators. These figures place the wide latest-period interval in context: even without a uniquely dominant review, nine clusters provide limited information. The two-way calculation clusters by security identifier, not ultimate issuer, and therefore does not absorb every possible cross-share-class or corporate-group dependence.

The deletion results are more weighting-sensitive. In 2012--2016, the event-weighted mean is -2.47\%, compared with an equal-review mean of -7.27\%; one review contains 41.3\% of the events. The jackknife standard error is 4.94 percentage points, compared with CR1's 3.50. The differing estimates describe different weighting targets and influential cluster composition. We do not choose the weight that yields the preferred economic narrative. Appendix~\ref{app:dependence} reports the later-period comparison in a compact table, and the audit files retain all periods.

\subsection{Returns through twenty trading days after effectiveness}

Table~\ref{tab:horizon} reports the new fixed-cohort extension. For additions in 2017--2026, the mean return before ED is 2.43\% and the non-overlapping ED-through-ED+20 component is 0.87\%. Their sum is 3.30\%, with a 95\% analytical interval of [1.31\%, 5.28\%] and an independently computed wild-bootstrap $p$-value below 0.001. The post-ED component itself has an interval of [-0.78\%, 2.51\%]. The cumulative result therefore does not establish a separately positive post-event drift.

\begin{table}[!htbp]
\centering
\caption{A fixed cohort decomposition through twenty days after implementation}\label{tab:horizon}
\small
\setlength{\tabcolsep}{5pt}
\begin{tabular}{llrrrrlr}
\toprule
Type & Period & $N$ & Before ED & After ED & Cumulative & CR1 interval & $p_{WCR}$ \\
\midrule
Addition & 2017--2026 & 169 & 2.43 & 0.87 & 3.30 & [1.31, 5.28] & $<0.001$ \\
Addition & 2022--2026 & 87 & 1.77 & -0.06 & 1.71 & [0.72, 2.69] & 0.010 \\
Deletion & 2017--2026 & 136 & -0.70 & -0.38 & -1.07 & [-3.68, 1.54] & 0.414 \\
Deletion & 2022--2026 & 75 & -0.93 & -0.95 & -1.88 & [-5.72, 1.96] & 0.373 \\
\bottomrule
\end{tabular}
\begin{minipage}{\linewidth}\footnotesize\vspace{5pt}
New exploratory audit. Before ED is CAR $[AD-1,ED-1]$; After ED is CAR $[ED,ED+20]$; Cumulative is their event-level sum, CAR $[AD-1,ED+20]$. Every entry in a row uses exactly the same complete-case events. All values are percentages. The interval is analytical CR1 with $t_{G-1}$; $p_{WCR}$ is a separately computed restricted wild cluster bootstrap $p$-value using 49,999 draws and the finite-simulation correction $(k+1)/(B+1)$. These are distinct inferential summaries. No claim of permanent or causal price impact follows from a nonzero finite-horizon cumulative return. Source: audit/results/paired\_horizon.csv.
\end{minipage}
\end{table}


The fixed cohort is essential to interpreting this decomposition. It ensures that the two components add at the event level and in their reported means. It does not turn the before-ED component into an estimate of a treatment effect, nor does it eliminate unrelated information arriving during the extended window. The economic content is that, for the measured later addition cohort and this benchmark, average cumulative performance has not returned to zero by the chosen endpoint. That statement concerns a finite horizon. It does not demonstrate a permanent price change.

The latest block illustrates how a longer window can change a statistical summary. Its cumulative mean is 1.71\%, with an analytical interval of [0.72\%, 2.69\%] and wild-bootstrap $p=0.010$. The post-ED mean is approximately -0.06\%. This estimate is nevertheless sensitive to the same observations that matter in the short-window analysis. Removing Sivers lowers the cumulative mean to 1.25\%, with an interval reaching below zero; removing MSCI overlaps gives 1.21\%, also with an interval spanning zero. The full latest-block estimate must not be promoted as a robust finding merely because one extended-window test is significant.

The later-decade cumulative mean is less sensitive in those particular checks. It is 3.07\% without Sivers, 2.51\% without MSCI overlaps and 3.14\% without the screened company-news observations; all three pointwise analytical intervals remain above zero. Appendix~\ref{app:horizon} gives the estimates and bootstrap tests. These exclusions were added after inspecting the extension and are explicitly exploratory. Positive results from several overlapping specifications are not several independent discoveries.

For deletions, the later-decade cumulative estimate is -1.07\%, with an interval of [-3.68\%, 1.54\%]. The latest-block value is -1.88\%, also imprecise. Hence the extension does not support a symmetric addition-deletion persistence claim. Nor can the imprecise deletion estimates rule out an economically relevant negative response. The strongest defensible addition to the thesis is the transparent horizon accounting and its sensitivity analysis, rather than a new claim that Swedish indexing has a permanent causal effect.

\subsection{What the auxiliary designs can resolve}

The dual-share-class exercise compares an indexed security with another class from the same firm. It can remove some firm-level information, but the comparison class may be affected by the event through substitution or shared investor attention. The available set contains 39 events from 13 firms. The primary sign-adjusted difference is 0.33 percentage points on an event-weighted basis and 0.68 on a firm-weighted basis; the respective placebo $p$-values are 0.407 and 0.322. Neither estimates a precisely zero effect. Excluding events with price-source changes leaves 29 events from nine firms and a mean of 0.04 percentage points.

The placebo construction draws alternative seasonal windows rather than exploiting random index assignment. Its event-wise draws do not fully preserve the joint history of repeated events within a firm. That limits the interpretation of its reference distribution. A short announcement-window difference of 0.70 percentage points has an unadjusted placebo $p$-value of 0.029, but the archived three-window adjustment raises it to 0.088. We do not elevate that isolated result above the primary window or use the imprecise primary contrast to conclude that demand has no effect.

The inherited share-issuance exercise and OMXS30 comparison are reported in Appendix~\ref{app:auxiliary}. They are informative about the available data but do not strengthen identification sufficiently to become main results. Share-count adjustments remain partly heuristic, and the small OMXS30 sample cannot establish equivalence to OMXSB. Keeping these exercises visible while limiting their argumentative weight avoids treating every available regression as corroboration.
\FloatBarrier
